% !TeX root = proceedings.tex
% Match the DOCX order and one-column title block followed by two-column text.
\vspace{.5em}
변지훈 · 유은빈 · 최서경 · 김준영 · 정소이*\par Jihun Byun · Eunbin Yu · Seokyung Choi · Junyoung Kim · Soyi Jung*\par
\vspace{.3em}
아주대학교 전자공학과\par Department of Electrical and Computer Engineering, Ajou University\par *Corresponding author: sjung@ajou.ac.kr\par

\twocolumn[{
\begin{@twocolumnfalse}
\centering
\vspace*{27bp}
{\proceedingstitlelatin\proceedingstitlehangul
 \fontsize{18bp}{28.08bp}\selectfont\bfseries
 4족 보행 로봇의 관절 토크 저하에 강건한\\
 교사, 학생 기반 강화학습 보행 제어\par}
\vspace{32bp}
{\proceedingsauthorlatin\proceedingsauthorhangul
 \fontsize{10bp}{17bp}\selectfont\bfseries
 \paperAuthorsKo\par}
{\fontsize{10bp}{17bp}\selectfont\bfseries
 \paperAffiliationKo\par}
{\fontsize{10bp}{18.12bp}\selectfont\bfseries
 \paperAuthorEmails\par}
\vspace{9bp}
{\fontsize{12bp}{18.72bp}\selectfont\bfseries
 Teacher–Student Reinforcement Learning for Robust Quadruped Locomotion under Joint Torque Degradation\par}
\vspace{19.8bp}
{\proceedingsauthorenglish
 \fontsize{10bp}{18.36bp}\selectfont\bfseries
 \paperAuthorsEn\par
 \paperAffiliationEn\par}
% Align the Abstract title and first right-column text at their prior midpoint.
\vspace{35.604bp}
\end{@twocolumnfalse}
}]
\vspace*{-16.204bp}
{\centering\fontsize{12bp}{18.72bp}\selectfont\bfseries Abstract\par}
\vspace{14.04bp}
\paperbodyfont
\indent% !TeX root = proceedings.tex
We propose a reinforcement learning locomotion policy that maintains balance, tracks commanded motion, and follows reference paths under nominal and single-joint torque degradation conditions. The policy combines current observations and history-based features through teacher–student joint training, without direct fault inputs. Policies trained with one seed each are compared in Isaac Gym using a Unitree A1 quadruped model on matched mixed-terrain tasks. Under full torque degradation, survival improved by 37.2 percentage points and mean pre-termination cross-track root mean square error decreased by 64.6\% relative to the nominally trained policy.
\par
% The proceedings template has no separate Korean summary or Keywords block.
% The original Korean summary and Keywords remain in frontmatter.tex for review.
